\subsection{Otra información relevante de la empresa donde se desarrolla el proyecto}

\subsubsection{Planificación y gestión de la entrevista}

Con el propósito de profundizar en el diagnóstico situacional, contrastar la perspectiva estratégica del negocio y sustentar las causas raíz de las deficiencias operativas en el cálculo de haberes, se diseñó e implementó una entrevista semiestructurada aplicada a los responsables directos de la gestión empresarial y financiera.

\vspace{0.3cm}
\noindent\textbf{Entrevistado:} Julio Flores Vicuña, \textbf{Gerente General de J\&P Periféricos S.A.C.}
\vspace{0.4cm}

\begin{enumerate}[leftmargin=*, label=\textbf{P\arabic*.}]
    \item \textbf{¿Cuáles son las principales dificultades de gestión interna que ha identificado en las empresas de su cartera de clientes?}
    
    \textit{Respuesta del Gerente General:} Los principales inconvenientes son los errores en cálculos de nómina, retrasos en pagos, registros imprecisos de horas/licencias y falta de integración tecnológica. Estos problemas generan reclamos, sanciones de SUNAFIL y pérdida de confianza en los equipos.

    \item \textbf{¿Qué motivos estratégicos llevaron a la empresa a optar por el desarrollo de una solución propia de planillas en lugar de contratar un software comercial existente?}
    
    \textit{Respuesta del Gerente General:} Un sistema de planillas personalizado permite a una empresa con muchos empleados reducir errores, cumplir con normativas locales (SUNAT/SUNAFIL), y adaptar cálculos a sus propias políticas de bonos, turnos y beneficios. Además, ofrece mayor seguridad de datos y flexibilidad frente a cambios legales o internos.

    \item \textbf{Desde su perspectiva como gerente, ¿cuáles representan los mayores desafíos al momento de supervisar el cumplimiento, asistencias e incidencias de los colaboradores?}
    
    \textit{Respuesta del Gerente General:} Tener la información actualizada en tiempo real de los integrantes de cada proceso y medir su desempeño.

    \item \textbf{¿Qué impacto operativo y económico espera que genere este sistema, tanto a nivel interno en J\&P Periféricos como para su futura comercialización?}
    
    \textit{Respuesta del Gerente General:} En el tema operativo mejoraríamos el orden de personal asignado a los procesos y en lo económico el recupero es a mediano plazo, recuperamos la inversión, con los costos que va a reducir el sistema.

    \item \textbf{¿Qué necesidades o problemáticas administrativas comunes comparten las empresas a las que estará orientada la venta del sistema?}
    
    \textit{Respuesta del Gerente General:} Manejo de personal es un problema general en las empresas que manejan proyectos de gran envergadura, cada empresa maneja su personal de acuerdo con los contratos; los contratos definen mucho la política de pago orientado en el sistema.
\end{enumerate}

\vspace{0.3cm}
\noindent\textbf{Entrevistado:} Juan Pablo Alcalá Juárez, \textbf{Contador de una de las empresas de la cartera de clientes de J\&P Periféricos S.A.C.}
\vspace{0.4cm}

\begin{enumerate}[leftmargin=*, label=\textbf{P\arabic*.}]
    \item \textbf{¿Cuánto tiempo promedio invierte mensualmente en el proceso completo de elaboración y validación de la planilla?}
    
    \textit{Respuesta del Contador:} Aproximadamente 5 horas, considerando la revisión de las incidencias de los colaboradores, el cálculo de las remuneraciones, la preparación de los archivos de pago, la generación y distribución de las boletas. 
    También considerando el tiempo de revisión de los cálculos y la verificación de la información antes de su envío a SUNAT y a los bancos.

    \item \textbf{¿Cuáles son las fases o variables (cálculo de horas extra, tardanzas, CTS, retenciones de ley) que demandan mayor tiempo y esfuerzo manual? }
    
    \textit{Respuesta del Contador:} La revisión de las incidencias y la verificación de los datos que se utilizarán para calcular la planilla. Es necesario comprobar las horas extras, tardanzas, 
    permisos e inasistencias, y verificar que correspondan al trabajador y al periodo adecuado. 
    Cuando se encuentran diferencias o información incompleta, se requiere consultar con Recursos Humanos y corregir los registros antes de continuar con los cálculos.
    
    \item \textbf{¿El proceso de elaboración de planilla lo realiza de manera individual o cuenta con el soporte de asistentes contables?}
    
    \textit{Respuesta del Contador:} Realizo el proceso con el apoyo de un asistente contable y del encargado de Recursos Humanos.

    \item \textbf{En caso de contar con apoyo, ¿cómo se distribuyen las tareas operativas (revisión de asistencia, cálculo de incidencias, digitación)? }
    
    \textit{Respuesta del Contador:} Al asistente contable le asigno la revisión de las incidencias y la digitación de los datos, 
    mientras que yo me encargo de realizar los cálculos y revisar los resultados. 
    Una vez revisada la planilla, le encargo al asistente la generación de los archivos de pago y las boletas.

    \item \textbf{¿Cuáles son los errores o inconsistencias más frecuentes al trabajar en hojas de cálculo y qué consecuencias conllevan al momento de efectuar los desembolsos? }
    
    \textit{Respuesta del Contador:} Los errores más frecuentes son el ingreso de datos incorrectos, la omisión de incidencias como tardanzas, permisos o inasistencias, y algunos casos de duplicación de registros. Estas inconsistencias pueden generar que nos demoremos en realizar los pagos, también la posibilidad de recibir sanciones por parte de SUNAFIL y pérdida de confianza por parte de los colaboradores.

\end{enumerate}